%% ma: L10-B14-0022
%% lop: 10
%% chuong: 5
%% bai: 14
%% dang: 10.14.2
%% loai: TN4
%% muc: TH
%% dap_an: "D"
%% nguon: "chuyen-de-cac-so-dac-trung-cua-mau-so-lieu-khong-ghep-nhom-toan-10-chuong-trinh-moi.pdf (GV. Phan Nhật Linh) (P04), Bài 03 Dạng 1 Phần I Câu 12"
%% kien_thuc: "Lớp 10 Bài 13 (tứ phân vị), Bài 14 (khoảng tứ phân vị)"
%% trang_thai: chinh-thuc
%% ghi_chu: "Tự duyệt theo yêu cầu giáo viên 22/10/2026."
\begin{ex}
Nhiệt độ ($^\circ$C) của một thành phố ghi nhận trong $10$ ngày qua lần lượt là:
\begin{center}
$24$\quad $21$\quad $30$\quad $34$\quad $28$\quad $35$\quad $33$\quad $36$\quad $25$\quad $27$
\end{center}
Khoảng tứ phân vị của mẫu số liệu bằng
\choice{$\Delta_Q=12$}{$\Delta_Q=11$}{$\Delta_Q=13$}{\True $\Delta_Q=9$}
\loigiai{Sắp xếp: $21$; $24$; $25$; $27$; $28$; $30$; $33$; $34$; $35$; $36$.

$Q_1=25$ (trung vị của $21$; $24$; $25$; $27$; $28$); $Q_3=34$ (trung vị của $30$; $33$; $34$; $35$; $36$). $\Delta_Q=34-25=9$. Chọn D.}
\end{ex}
